% Generated draft: requires mathematical and source review.
\documentclass[11pt]{article}
\usepackage[margin=25mm]{geometry}
\usepackage{fontspec}
\setmainfont{Latin Modern Roman}
\usepackage{amsmath,amssymb,mathtools,unicode-math}
\setmathfont{Latin Modern Math}
\usepackage{enumitem,needspace,xurl,hyperref}
\setlength{\parindent}{0pt}
\setlength{\parskip}{5pt}
\setlength{\emergencystretch}{3em}
\newcommand{\N}{\mathbb{N}}
\newcommand{\Z}{\mathbb{Z}}
\newcommand{\Q}{\mathbb{Q}}
\newcommand{\R}{\mathbb{R}}
\newcommand{\C}{\mathbb{C}}
\newcommand{\F}{\mathbb{F}}
\newcommand{\A}{\mathbb{A}}
\newcommand{\PP}{\mathbb{P}}
\newcommand{\E}{\mathbb{E}}
\newcommand{\eps}{\varepsilon}
\newcommand{\dd}{\,\mathrm{d}}
\begin{document}
\section*{UnsolvedMath ID: 30006211}
\textbf{Source problem code:} catalogue entry
\subsection{Definitions and mathematical statement}
For an integer \(n\geq 1\), a polynomial \(h\in\R[x_1,\ldots,x_n]\) is multi-affine if its degree in each variable is at most one. Equivalently, it has the form
\[
h(x)=\sum_{S\subseteq\{1,\ldots,n\}}a_S\prod_{i\in S}x_i,\qquad a_S\in\R.
\]
Here the empty product is \(1\). Interpret a degree bound \(d\) as a total-degree bound: \(a_S=0\) whenever \(|S|>d\). This formulation also includes constant and zero polynomials without needing to assign a degree to the zero polynomial.

For two such polynomials define their common real zero set by
\[
Z_{\R}(f,g)=\{x\in\R^n:f(x)=g(x)=0\}.
\]
Let \(b_0(Z_{\R}(f,g))\) denote its number of connected components in the Euclidean topology, with \(b_0(\varnothing)=0\).

The question is whether
\[
\begin{gathered}
\forall d\in\Z_{\geq0}\ \exists B_d\in\Z_{\geq0}\ \forall n\in\Z_{\geq1}\ \\
\forall f,g\in\R[x_1,\ldots,x_n]\text{ multi-affine of total degree at most }d,\\
b_0\bigl(Z_{\R}(f,g)\bigr)\leq B_d.
\end{gathered}
\]
The bound must be independent of \(n\) and of the polynomial coefficients. A bound in terms of a common degree bound is equivalent to allowing dependence on the two individual degree bounds, by taking their maximum.

The supplied excerpt does not explicitly specify the degree convention or coefficient field; the displayed formulation uses total degree and real coefficients as the conventional interpretation of the stated real-zero-set question. These interpretations require confirmation against the report. No smoothness, independence, or nondegeneracy hypotheses are supplied.
\subsection{Short English statement}
Can the number of connected components of the common zero set in \(\R^n\) of two real multi-affine polynomials of bounded total degree be bounded independently of the number of variables? The supplied record labels this question open; this statement-only draft requires human mathematical review.
\subsection{Sources}
\begin{itemize}
\item Local UnsolvedMath registry: supplied complete record with permanent ID \(30006211\), title \emph{Real Components of Two Multi-Affine Polynomial Zero Sets}, source problem code \texttt{catalogue entry}, in the Oberwolfach Reports Open Problems collection. The record gives proposed year 2025 and supplies no individual proposer.
\item Supplied authoritative source citation: \emph{MFO-RIMS Tandem Workshop: Optimization, Theoretical Computer Science and Algebraic Geometry: Convexity and Beyond} (2025), Oberwolfach Reports, pp.\ 405--450, DOI \texttt{10.4171/owr/2025/9}. Source URL and literature-source link labelled OWR report: \url{https://doi.org/10.4171/owr/2025/9}. The supplied background describes these reports as documenting mathematical workshops and their open-problem discussions. No precise page within the report is supplied for this question.
\item Supplied literature-source link labelled Basu--Perrucci: \url{https://arxiv.org/abs/2204.01595}. The record supplies no title or precise result locator for this reference.
\item Literature assessment checked at \texttt{2026-08-22}: the supplied curation labels the problem open and reports that the one-polynomial bound and unbounded three-polynomial constructions are established, whereas the dimension-independent bound for two multi-affine polynomials remains open. This is supplied dated curation derived from the OWR material, not an independently verified literature assessment. The record explicitly calls for independent verification of the time-sensitive status and cited claims.
\end{itemize}
\end{document}
